\documentclass[12pt]{article}

% Documento específico para a entrega do Seminário de Andamento.
% As seções acadêmicas são compartilhadas com o artigo completo para evitar
% divergência de conteúdo entre as duas versões.
\usepackage{sbc-template}
\usepackage[T1]{fontenc}
\usepackage[utf8]{inputenc}
\usepackage[brazil]{babel}
\usepackage{graphicx}
\usepackage{url}
\usepackage{xurl}
\usepackage{xcolor}
\usepackage{booktabs}
\usepackage{tabularx}
\usepackage{array}
\usepackage{longtable}
\usepackage{float}
\usepackage{amsmath}
\usepackage{enumitem}
\usepackage{microtype}
\usepackage{tikz}
\usetikzlibrary{arrows.meta,positioning,shapes.geometric,fit}
\usepackage[hidelinks]{hyperref}

% Centralize aqui os dados de identificacao antes da entrega.
% Nao remova um marcador PENDENTE sem confirmar a informacao.
\newcommand{\TituloTCC}{Prevenção de abastecimento incompatível por meio de consulta veicular e intertravamento de bicos: projeto e avaliação de um protótipo}
\newcommand{\AutorTCC}{Gabriel Nahmur}
\newcommand{\EmailAutor}{\texttt{[PENDENTE: e-mail acadêmico]}}
\newcommand{\OrientadorTCC}{\texttt{[PENDENTE: nome completo do(a) orientador(a)]}}
\newcommand{\InstituicaoTCC}{Instituto Federal de Educação, Ciência e Tecnologia Sul-rio-grandense (IFSul)}
\newcommand{\CampusTCC}{Câmpus Bagé}
\newcommand{\CursoTCC}{Curso Superior de Tecnologia em Análise e Desenvolvimento de Sistemas}
\newcommand{\CidadeDataTCC}{Bagé, 2026}


\hypersetup{
  pdftitle={Seminário de Andamento: \TituloTCC},
  pdfauthor={\AutorTCC},
  pdfsubject={Seminário de Andamento - recorte acadêmico para entrega},
  pdfkeywords={abastecimento incompatível; intertravamento; consulta veicular; Design Science Research; análise e projeto de software}
}

\addto\captionsbrazil{\renewcommand{\refname}{Referências}}
\renewcommand{\arraystretch}{1.16}
\newcolumntype{L}[1]{>{\raggedright\arraybackslash}p{#1}}
\newcolumntype{Y}{>{\raggedright\arraybackslash}X}
\setlist[itemize]{leftmargin=1.27cm}
\setlist[enumerate]{leftmargin=1.27cm}
\emergencystretch=2em

\newcommand{\pendente}[1]{\textcolor{red!65!black}{\textbf{[PENDENTE: #1]}}}
\newcommand{\naorealizado}{\textit{Não realizado até esta versão do seminário.}}
\newcommand{\fontepropria}{Fonte: elaborado pelo autor (2026).}

\title{Seminário de Andamento\\[0.35em]\large \TituloTCC}
\author{\AutorTCC}
\address{\InstituicaoTCC\\
  \CampusTCC\\
  \CursoTCC\\
  Orientação: \OrientadorTCC\\
  \CidadeDataTCC
  \email{\EmailAutor}
}

\begin{document}
\raggedbottom

\maketitle

\begin{center}
\fcolorbox{black!55}{black!3}{%
  \parbox{0.88\textwidth}{\small\textbf{Documento parcial para o Seminário de Andamento.}
  Esta versão contém somente contextualização e objetivos, fundamentação
  teórica, revisão de literatura, metodologia parcial, documentos de análise e
  projeto do software e cronograma. Campos em vermelho permanecem pendentes de
  confirmação formal.}}
\end{center}

\begin{resumo}
O abastecimento de um veículo com combustível incompatível pode gerar
transtornos operacionais, custos de reparação e impacto no atendimento ao
cliente. Este trabalho propõe um protótipo seguro de bancada que combina
consulta de dados veiculares, política de compatibilidade e intertravamento
físico independente no suporte dos bicos. O artefato não comanda bombeamento
nem modifica a automação da empresa. A pesquisa adota \textit{Design Science
Research}, revisão de literatura e estudo de caso em uma rede de postos
anonimizada. Esta versão de andamento apresenta a contextualização, os
objetivos, a fundamentação, o protocolo da revisão, a metodologia parcial, os
requisitos, a arquitetura, a modelagem e o cronograma. Dados históricos,
entrevistas e ensaios físicos dependem de autorizações e etapas posteriores;
portanto, nenhuma hipótese é considerada concluída neste documento.

\par\smallskip\noindent\textbf{Palavras-chave:} Abastecimento incompatível,
Intertravamento, Consulta veicular, Análise e projeto de software, Design
Science Research.
\end{resumo}

% (1) Contextualização e (2) definição de objetivos.
\section{Introdução}
\label{sec:introducao}

O abastecimento é uma atividade cotidiana, mas envolve uma decisão que precisa
ser correta antes do início da transferência do combustível: associar o produto
selecionado à configuração admitida pelo veículo. Quando essa associação falha,
o procedimento posterior depende do tipo e da quantidade de combustível, do
funcionamento do motor e das recomendações do fabricante. A literatura registra
que o erro não pode ser explicado apenas por falta de atenção. Em uma análise
de tarefa e entrevistas com 23 pessoas que haviam abastecido veículos leves com
o produto errado no Reino Unido, \cite{adamsDavid2007} identificaram combinações
de fatores humanos e ambientais capazes de formar caminhos para o erro. Os
autores defenderam intervenções no leiaute e no equipamento para interromper
esses caminhos.

O presente trabalho nasce de um problema relatado em uma rede de postos de
combustíveis do sul do Brasil: durante o atendimento, pode ocorrer a seleção de
um produto incompatível com o veículo. A existência, frequência e composição
dos custos locais ainda serão verificadas por registros administrativos; por
isso, esta versão não apresenta número de ocorrências nem valor de prejuízo.
Essa separação é importante, pois dados internacionais ou relatos isolados não
podem ser convertidos diretamente em prevalência para a organização estudada.

Como possível barreira anterior à liberação, propõe-se consultar a placa do
veículo em um serviço autorizado, transformar o código retornado em uma política
explícita de compatibilidade e liberar temporariamente apenas o bico solicitado
e compatível. O serviço Consulta Online SENATRAN disponibiliza código e descrição
do combustível na modalidade detalhada \cite{serproDadosVeiculo2026}. Entretanto,
seu uso não se resume a uma chamada técnica: requer caso de uso aprovado,
contratação do Serpro e de uma Gerenciadora de Consentimento e Ciência (GCC) e
atendimento às condições de governança aplicáveis \cite{senatran1392025,serproComoContratar2026}.

O artefato acadêmico será inicialmente construído em bancada seca, sem
combustível e em baixa tensão. A trava física representa o intertravamento, mas
não constitui solução pronta para instalação em uma bomba. Uma implantação real
exigiria integração com o fabricante ou controlador de pista, classificação do
local e equipamentos adequados à atmosfera potencialmente explosiva
\cite{inmetro1152022,nr20}.

\subsection{Problema e pergunta de pesquisa}

O problema de pesquisa consiste na ausência de uma barreira automatizada que,
no contexto estudado, verifique a compatibilidade antes de autorizar o acesso ao
bico. A investigação é orientada pela seguinte pergunta:

\begin{quote}
\textit{Como um sistema baseado em consulta veicular e intertravamento pode
impedir a liberação de combustíveis incompatíveis em ambiente controlado, e em
quais condições sua adoção apresenta viabilidade econômica para uma rede de
postos?}
\end{quote}

Foram formuladas duas hipóteses verificáveis:

\begin{itemize}
  \item \textbf{H1:} na matriz de testes de bancada, o sistema impedirá todas as
  liberações incompatíveis e não liberará automaticamente nenhum bico diante de
  resposta inválida, ambígua ou indisponível;
  \item \textbf{H2:} em pelo menos um dos cenários de sensibilidade, as perdas
  evitáveis estimadas superarão o custo anualizado da solução.
\end{itemize}

H1 somente será aceita se não houver liberação indevida nos casos definidos
antes dos ensaios. H2 será avaliada por cenários, e não como economia observada.
Ambas permanecem sem teste conclusivo nesta etapa.

\subsection{Objetivos}

O objetivo geral é projetar, implementar e avaliar um protótipo para prevenção
de abastecimento incompatível, integrando consulta veicular, controle de
autorização e análise econômica baseada em dados históricos anonimizados.

São objetivos específicos:

\begin{enumerate}
  \item caracterizar o processo de abastecimento e, quando existirem registros
  suficientes, as ocorrências e os custos documentados no estudo de caso;
  \item mapear evidências científicas e documentação técnica sobre erro de
  abastecimento, fatores humanos, identificação veicular e intertravamento;
  \item elicitar requisitos funcionais, não funcionais, de privacidade e de
  segurança para o artefato;
  \item desenvolver um simulador do provedor veicular e um protótipo em bancada
  com política de falha para o estado bloqueado;
  \item avaliar o comportamento funcional, o tratamento de falhas, a latência,
  a usabilidade e a repetibilidade do protótipo; e
  \item estimar custos e ponto de equilíbrio sob diferentes percentuais de
  prevenção, distinguindo gastos comprovados de premissas.
\end{enumerate}

\subsection{Justificativa e contribuições esperadas}

A relevância prática decorre do potencial de agir antes do dano, quando a
escolha ainda pode ser bloqueada. A relevância acadêmica está na construção e
avaliação explícita de um artefato sociotécnico: não basta reconhecer a placa;
é necessário tratar códigos desconhecidos, indisponibilidade, exceções
operacionais, dados pessoais e o limite entre protótipo e equipamento de campo.

As contribuições esperadas são: (i) uma caracterização documentada do problema
na empresa; (ii) um protocolo e uma síntese da literatura; (iii) requisitos e
arquitetura auditáveis; (iv) um protótipo reprodutível de bancada; (v) uma matriz
de testes com estados de falha; e (vi) um modelo econômico transparente. Essas
contribuições são metas do trabalho e não resultados já alcançados.


% (3) Referencial teórico.
\section{Fundamentação teórica}
\label{sec:fundamentacao}

\subsection{Erro humano e barreiras de prevenção}

\cite{reason1990} diferencia formas de falha associadas à execução e ao
planejamento e mostra que o desempenho humano é influenciado pelo sistema em
que a tarefa ocorre. Essa perspectiva evita reduzir o problema a uma culpa
individual e direciona a análise para sinais, sequência da tarefa, interface,
equipamento e condições de trabalho. No abastecimento, \cite{adamsDavid2007}
aplicaram análise hierárquica da tarefa e encontraram fatores físicos e
psicológicos que, isolados ou combinados, abrem caminhos para a seleção errada.

O intertravamento proposto é uma barreira de prevenção: a ação física depende de
uma autorização calculada antes da retirada do bico de seu suporte. A barreira
atua sobre o acesso ao bico, não sobre o bombeamento ou a vazão. Isso não elimina
treinamento, conferência e procedimentos. Também pode criar novos modos de
falha, como indisponibilidade do serviço, código veicular não reconhecido,
comando atrasado ou uso abusivo de uma exceção. Por isso, o projeto adota os
princípios de autorização explícita, menor privilégio e padrão seguro diante de
falhas, coerentes com princípios clássicos de proteção de sistemas
\cite{saltzerSchroeder1975}.

\subsection{Consulta veicular e política de compatibilidade}

O serviço Consulta Online SENATRAN expõe operações de consulta a bases mantidas pela
Secretaria Nacional de Trânsito. No inventário público, os campos ``Código
Combustível'' e ``Combustível Descrição'' aparecem na consulta detalhada de
veículo \cite{serproDadosVeiculo2026}. O retorno do provedor, contudo, não deve
acionar diretamente um dispositivo. Entre a consulta e o atuador haverá uma
camada de política que:

\begin{enumerate}
  \item valida formato, atualidade e completude da resposta;
  \item converte códigos externos para categorias internas versionadas;
  \item representa veículos com um ou mais combustíveis admitidos;
  \item compara o produto solicitado à lista de compatibilidade; e
  \item produz uma decisão auditável: autorizado, negado ou revisão necessária.
\end{enumerate}

Essa separação permite substituir o provedor real por dados sintéticos durante
o desenvolvimento e testar códigos novos sem expor dados pessoais. Também evita
acoplar a segurança física a detalhes instáveis da API.

\subsection{Sistemas ciberfísicos, mensageria e estado seguro}

O protótipo combina software de negócio, comunicação e atuadores e foi concebido
para que uma decisão digital possa comandar um efeito físico observável. A
validação do movimento mecânico, porém, permanece pendente. A comunicação entre o backend
e o controlador será representada por mensagens MQTT, protocolo padronizado pela
OASIS com níveis de qualidade de serviço e sessões próprias
\cite{oasisMqtt2019}. Esses recursos não garantem, sozinhos, que um comando seja
executado uma única vez. O projeto acrescenta identificador idempotente,
validade temporal, confirmação e bloqueio local por expiração.

O estado seguro comandado do protótipo é manter desenergizadas as saídas de
abertura de todas as travas. Inicialização, perda de conexão, resposta ambígua,
comando vencido e término da sessão conduzem a esse estado. Se um bico já
estiver fora do suporte, desenergizar a saída não o recoloca fisicamente; sua
devolução continua sendo uma ação humana. Uma liberação excepcional não altera
essa premissa: ela exige supervisor autenticado, justificativa e trilha de
auditoria.

\subsection{Proteção de dados e acesso regulado}

A placa pode viabilizar associação a uma pessoa natural no contexto do
tratamento e, por prudência, será tratada como dado pessoal. A Lei Geral de
Proteção de Dados Pessoais estabelece, entre outros, os princípios de finalidade,
adequação, necessidade, segurança, prevenção e responsabilização
\cite{lgpd2018}. A Portaria SENATRAN n\textsuperscript{o} 139/2025 disciplina o acesso aos sistemas
e subsistemas, incluindo avaliação do caso de uso, controles de segurança e
atuação do Serpro e de uma GCC \cite{senatran1392025}.

Consequentemente, a versão de bancada utiliza placas e respostas sintéticas.
Na integração autorizada, a placa será usada apenas durante a decisão e não
constará do log padrão. Também não serão persistidos resposta bruta, nome de
proprietário, Renavam ou chassi. A base legal, os avisos ao titular, a retenção e
as responsabilidades deverão refletir o caso de uso efetivamente aprovado, não
uma suposição do projeto.

\subsection{Segurança da instalação}

A NR-20 define requisitos de segurança e saúde no trabalho com inflamáveis e
combustíveis e contempla o conceito de área classificada \cite{nr20}. A Portaria
Inmetro n\textsuperscript{o} 115/2022 trata da avaliação da conformidade de equipamentos elétricos
para atmosferas explosivas \cite{inmetro1152022}. Assim, ESP32, servo, solenoide
e tablet comuns não serão instalados na zona de abastecimento. O ensaio acadêmico
ocorrerá sem líquido inflamável, e uma eventual evolução de campo dependerá de
projeto específico do intertravamento independente, profissionais habilitados e
equipamentos certificados para a classificação do local. Essa evolução não
pressupõe alteração da bomba ou integração com a automação da empresa.


% (4) Revisão de literatura.
\section{Revisão de literatura e trabalhos relacionados}
\label{sec:revisao}

\subsection{Protocolo de busca e seleção}

Será conduzido um mapeamento da literatura conforme as recomendações atualizadas
de \cite{petersenEtAl2015}. O protocolo será congelado antes da seleção final e
registrará bases, strings, filtros, datas, duplicatas, decisões por título e
resumo, leitura completa, avaliação de qualidade e extração. As bases definidas
são IEEE Xplore, ACM Digital Library, Scopus e SciELO. Google Scholar será usado
apenas de forma complementar e para encadeamento de referências.

As strings combinam três grupos: (i) \textit{misfueling}, \textit{wrong fuel} e
\textit{fueling error}; (ii) prevenção, intertravamento, autorização, ergonomia
e fatores humanos; e (iii) identificação veicular, placa, RFID e dispensador.
Serão incluídos estudos revisados por pares, em português, inglês ou espanhol,
com texto completo e relação direta com ao menos uma questão da revisão.
Patentes, publicidade, notícias e manuais serão excluídos da síntese científica,
mas poderão ser analisados separadamente como documentação técnica.

Cada estudo receberá avaliação qualitativa quanto à clareza do objetivo,
descrição de contexto e amostra, reprodutibilidade do método, coerência das
métricas, separação entre resultados e interpretação e discussão de limitações.
O protocolo completo está versionado em
\texttt{research/protocol/revisao-literatura.md}; o diário e a matriz de extração
ficam em \texttt{research/literature/}.

\pendente{executar as buscas em todas as bases e inserir o fluxograma de seleção
com contagens verificáveis.}

\subsection{Síntese preliminar}

A Tabela~\ref{tab:trabalhos-relacionados} reúne as fontes centrais já
verificadas. Ela não representa o resultado final do mapeamento.

\begin{table}[H]
  \centering
  \caption{Trabalhos e documentos relacionados já verificados}
  \label{tab:trabalhos-relacionados}
  \small
  \begin{tabularx}{\textwidth}{@{}L{3.2cm}L{3.3cm}Y@{}}
    \toprule
    \textbf{Fonte} & \textbf{Método ou artefato} & \textbf{Relação com esta pesquisa} \\
    \midrule
    Adams e David (2007) & Análise hierárquica da tarefa e dados de 23 pessoas que cometeram o erro & Evidencia a combinação de fatores humanos e ambientais e recomenda modificar leiaute ou equipamento. \\
    Mikkonen e Engman (2013) & Análise técnica das consequências da troca entre gasolina e diesel & Sustenta que pequenas quantidades de gasolina podem degradar propriedades relevantes do diesel e que a gravidade depende do cenário. \\
    Taylor (1978) & Levantamento em postos de grandes cidades canadenses & Oferece registro histórico de incidência e de alterações em bocais, mas não permite inferir frequência atual ou brasileira. \\
    Gilbarco (manual técnico) & Identificação automática por RFID e autorização da bomba & Documenta uma abordagem comercial distinta. Não é evidência científica de eficácia nem define a arquitetura deste projeto, que não comanda bombeamento. \\
    \bottomrule
  \end{tabularx}
  \\[2pt]\fontepropria
\end{table}

O estudo de \cite{adamsDavid2007} é o mais diretamente relacionado à origem do
problema: trata o abastecimento como uma sequência suscetível a condições que
favorecem o erro. Já \cite{mikkonenEngman2013} concentra-se nas consequências
técnicas da contaminação entre gasolina e diesel, o que ajuda a justificar a
atuação preventiva, mas não avalia uma barreira automatizada. A documentação
\cite{gilbarcoFuelPoint2018} mostra que controladores comerciais já realizam
autorização eletrônica de abastecimento por identificação automática. A solução
proposta neste TCC difere ao investigar consulta veicular por placa no contexto
regulatório brasileiro, política explícita de compatibilidade, comportamento
\textit{fail-closed} e análise econômica para uma rede específica.

Até esta etapa, não foi identificado estudo revisado por pares que avalie
conjuntamente esses quatro elementos. Essa observação é um \textbf{achado
preliminar}, limitado às fontes-semente, e somente poderá ser apresentada como
lacuna consolidada após a execução integral do protocolo.


% (5) Metodologia parcial e documentos de análise/projeto do software.
\section{Metodologia}
\label{sec:metodologia}

\subsection{Delineamento da pesquisa}

A pesquisa é aplicada, de abordagem mista e objetivo exploratório-descritivo.
O eixo de construção seguirá a \textit{Design Science Research Methodology}
(DSRM), que organiza o trabalho em identificação do problema, definição dos
objetivos, projeto e desenvolvimento, demonstração, avaliação e comunicação
\cite{peffersEtAl2007}. A adoção da DSRM é adequada porque a contribuição inclui
um artefato destinado a resolver um problema organizacional, cuja utilidade
precisa ser demonstrada e avaliada \cite{hevnerEtAl2004}.

Em paralelo, será realizado um estudo de caso único, com unidades de análise
incorporadas correspondentes a funções e, quando autorizado, a diferentes
unidades da rede. O estudo de caso contextualizará o processo, os incidentes e
as restrições organizacionais, seguindo recomendações de planejamento,
triangulação e relato de \cite{runesonHost2009}. A empresa será identificada no
texto apenas como ``uma rede de postos de combustíveis do sul do Brasil''.

\begin{table}[H]
  \centering
  \caption{Etapas da DSRM e evidências previstas}
  \label{tab:dsrm}
  \small
  \begin{tabularx}{\textwidth}{@{}L{3.1cm}YL{2.5cm}@{}}
    \toprule
    \textbf{Etapa} & \textbf{Evidência ou produto} & \textbf{Estado} \\
    \midrule
    Identificar problema & Registros anonimizados, entrevistas e descrição do processo & Parcial; dados pendentes \\
    Definir objetivos & Pergunta, hipóteses, objetivos e critérios de aceitação & Elaborado \\
    Projetar e desenvolver & Requisitos, arquitetura, simuladores, software e bancada & Software parcial; bancada pendente \\
    Demonstrar & Roteiro reproduzível com casos compatíveis, incompatíveis e falhas & Componentes simulados; roteiro integral pendente \\
    Avaliar & Testes funcionais, falhas, latência, usabilidade e cenários econômicos & Verificação automatizada parcial \\
    Comunicar & Artigo, repositório, seminário e defesa & Em andamento \\
    \bottomrule
  \end{tabularx}
  \\[2pt]\fontepropria
\end{table}

\subsection{Caracterização do processo e entrevistas}

Serão convidados, conforme autorização e disponibilidade, frentistas de mais de
uma unidade, gestores operacionais, representante da direção e profissional do
setor financeiro. A amostra será intencional e buscará diversidade de papéis,
turnos e unidades, sem pretensão de representatividade estatística. O roteiro
semiestruturado abrangerá sequência de atendimento, pontos de conferência,
causas percebidas, tratamento de exceções, consequências e aceitação de uma
barreira automatizada.

As entrevistas somente ocorrerão após autorização escrita da empresa,
orientação formal sobre o enquadramento ético e, quando aplicável, aprovação do
Comitê de Ética em Pesquisa (CEP). Serão adotados consentimento, participação
voluntária e plano de anonimização, retenção e descarte coerentes com a
Resolução CNS n\textsuperscript{o} 510/2016 \cite{cns5102016}. Gravação será opcional e dependerá
de consentimento específico. A análise utilizará codificação temática; as
categorias, decisões de codificação e trechos anonimizados serão registrados
para auditabilidade, sem expor pessoas ou unidades.

\pendente{registrar a decisão do orientador/CEP, a autorização da empresa, os
critérios finais de recrutamento e o número efetivamente entrevistado.}

\subsection{Dados históricos e custos}

A coleta documental solicitará um conjunto mínimo: identificador artificial da
ocorrência, período, unidade pseudonimizada, categorias de combustível,
momento de detecção, fonte documental e componentes comprovados do custo. Placa,
nome de cliente, nome de funcionário, chassi, Renavam e relato livre
identificável não serão inseridos na base de pesquisa.

Os valores serão classificados em custo documentado, estimativa e impacto
qualitativo. Somente itens suportados por registros serão somados como perdas
realizadas. Quando a série abranger datas diferentes, a correção para a data-base
usará o IPCA oficial \cite{ibgeIpca2026}, conforme a Equação~\ref{eq:ipca}:

\begin{equation}
  V_{corrigido}=V_{evento}\times
  \frac{IPCA_{data\mbox{-}base}}{IPCA_{evento}}.
  \label{eq:ipca}
\end{equation}

A taxa de incidentes por abastecimento somente será calculada se numerador e
denominador tiverem mesma população, período e unidades. Sem essa compatibilidade,
serão apresentadas apenas contagens e custos absolutos. O dicionário e as regras
de qualidade estão em \texttt{research/data/dicionario-dados.md}.

\subsection{Construção do artefato}

O desenvolvimento é incremental. Já foram implementados o domínio, o provedor
veicular simulado, a política de compatibilidade, o simulador de bicos, a API,
a auditoria minimizada em SQLite e uma PWA. Também existe firmware experimental,
cuja lógica portátil foi verificada em computador e cujo alvo \path{esp32dev}
foi compilado com PlatformIO. Autenticação individual,
persistência da versão da política, gateway MQTT no backend e integração entre
os componentes permanecem pendentes. A bancada física somente será montada após
validação do circuito e dos componentes com o orientador. A integração real com
SENATRAN permanecerá desativada até a obtenção das autorizações e credenciais
necessárias.

Cada cenário de demonstração terá entrada, estado inicial, eventos, decisão
esperada, comando esperado e estado final. A política e os contratos de
interface serão mantidos independentes da tecnologia física, permitindo trocar
um atuador simulado por um controlador de bancada sem modificar a regra de
negócio.

\subsection{Plano de avaliação}

A avaliação combinará cinco dimensões:

\begin{enumerate}
  \item \textbf{correção funcional:} decisão para veículos e produtos
  compatíveis, incompatíveis, multicombustível e desconhecidos;
  \item \textbf{segurança diante de falhas:} timeout, resposta inválida, perda
  de comunicação, reinicialização, mensagem duplicada e expiração;
  \item \textbf{desempenho:} latência ponta a ponta com mediana, percentil 95 e
  máximo, separando provedor, backend e dispositivo;
  \item \textbf{usabilidade:} conclusão da tarefa, erros, tempo e percepção dos
  participantes autorizados, com estatística descritiva; e
  \item \textbf{viabilidade econômica:} custo anualizado e perdas evitáveis sob
  cenários de 25\%, 50\%, 75\% e 100\% de prevenção.
\end{enumerate}

Sejam $Q$ o número anual de consultas, $c_q$ o custo unitário, $C_{GCC}$ a
remuneração anual da GCC, $C_{infra}$ infraestrutura, $C_{man}$ manutenção e
$C_{hw,a}$ a parcela anualizada do hardware. O custo anual da solução será:

\begin{equation}
  C_{\mathrm{sol}}=Q c_q+C_{GCC}+C_{infra}+C_{man}+C_{hw,a}.
  \label{eq:custo-solucao}
\end{equation}

Para uma perda histórica anual corrigida $L$ e percentual de prevenção $p$, o
benefício líquido contrafactual será $B(p)=pL-C_{\mathrm{sol}}$. Quando $L>0$, o ponto
de equilíbrio será $p^{*}=C_{\mathrm{sol}}/L$. O modelo não incluirá ganho de imagem,
fidelização ou risco jurídico sem método e evidência próprios.

\subsection{Ameaças à validade}

As principais ameaças previstas são registros incompletos, memória seletiva nas
entrevistas, pequena amostra, influência do pesquisador que também trabalha na
empresa, uso de simulador, diferença entre bancada e pista e mudanças de preço
ou regulação. As estratégias de mitigação incluem triangulação entre fontes,
roteiro comum, registro de decisões, dados sintéticos versionados, repetição dos
ensaios, relato de casos divergentes, revisão pelo orientador e atualização das
fontes oficiais na data da análise. Nenhuma generalização para outros postos ou
uma instalação real será feita sem evidência correspondente.

\section{Requisitos}
\label{sec:requisitos}

As Tabelas~\ref{tab:rf} e~\ref{tab:rnf} especificam requisitos-alvo do
artefato completo; sua presença não significa que todos estejam implementados
na versão de andamento. Os atores preliminares são: \textbf{frentista}, que solicita autorização;
\textbf{supervisor}, que trata exceções justificadas; \textbf{administrador},
que mantém usuários, identificadores dos conjuntos de bicos e política;
\textbf{auditor/pesquisador}, que acessa
dados minimizados; \textbf{provedor veicular}, simulado ou autorizado; e
\textbf{controlador de bancada}, que executa comandos restritos. Cliente e
titular de dados não são usuários do protótipo, mas são partes interessadas em
segurança, transparência e proteção de dados.

\begin{table}[H]
  \centering
  \caption{Requisitos funcionais preliminares}
  \label{tab:rf}
  \small
  \begin{tabularx}{\textwidth}{@{}L{1.1cm}Y@{}}
    \toprule
    \textbf{ID} & \textbf{Requisito} \\
    \midrule
    RF01 & Receber identificador idempotente da solicitação, placa, combustível, frentista e conjunto de bicos (\path{pump_id}), validando os campos. \\
    RF02 & Consultar um \textit{VehicleDataProvider} sem expor credenciais ao cliente web. \\
    RF03 & Converter códigos externos em categorias internas versionadas e registrar a versão da regra. \\
    RF04 & Retornar \path{AUTHORIZED}, \path{DENIED} ou \path{REVIEW_REQUIRED} com motivo padronizado. \\
    RF05 & Autorizar temporariamente a abertura somente da trava do suporte associado ao produto solicitado e compatível. \\
    RF06 & Desenergizar todas as saídas de abertura após expiração, encerramento, reinicialização ou perda de comunicação. \\
    RF07 & Permitir override somente a supervisor autenticado, com justificativa e validade limitada. \\
    RF08 & Rejeitar comando duplicado ou vencido e registrar a confirmação do dispositivo. \\
    RF09 & Manter trilha de auditoria minimizada, sem placa ou resposta veicular bruta no log padrão. \\
    RF10 & Operar com provedores e gateways simulados para desenvolvimento e demonstração sem dado real. \\
    \bottomrule
  \end{tabularx}
  \\[2pt]\fontepropria
\end{table}

\begin{table}[H]
  \centering
  \caption{Requisitos não funcionais e restrições preliminares}
  \label{tab:rnf}
  \small
  \begin{tabularx}{\textwidth}{@{}L{1.2cm}L{2.9cm}Y@{}}
    \toprule
    \textbf{ID} & \textbf{Atributo} & \textbf{Critério} \\
    \midrule
    RNF01 & Segurança funcional & Estado inicial e estado após falha devem ser \path{LOCK_ALL}. \\
    RNF02 & Privacidade & Usar minimização e não persistir placa, chassi, Renavam, proprietário ou resposta bruta por padrão. \\
    RNF03 & Segurança lógica & Credenciais externas ao repositório, autenticação individual e menor privilégio. \\
    RNF04 & Auditabilidade & Toda decisão e override deve registrar horário, ator, motivo, versão da política e resultado do dispositivo. \\
    RNF05 & Interoperabilidade & Provedor e gateway devem respeitar interfaces substituíveis, sem acoplamento do domínio à SENATRAN ou ao ESP32. \\
    RNF06 & Desempenho & Medir mediana, P95 e máximo; o limite operacional será definido após linha de base, sem número arbitrário nesta versão. \\
    RNF07 & Reprodutibilidade & Casos sintéticos, configuração de bancada e testes devem ser versionados. \\
    RNF08 & Segurança física & Operar apenas em bancada seca e de baixa tensão; não comandar bombeamento nem automação de pista; instalação real fica fora do escopo. \\
    \bottomrule
  \end{tabularx}
  \\[2pt]\fontepropria
\end{table}

Na versão verificada em 31 de julho de 2026, o backend já executa decisões com
provedor e intertravamento simulados, rejeita comandos duplicados ou vencidos e
mantém auditoria minimizada. A repetição de uma solicitação com o mesmo
\path{client_request_id} e o mesmo conteúdo recupera a decisão existente sem
novo acionamento; conteúdo divergente gera conflito. Entretanto, RF03, RF07,
RF08, RNF03 e RNF04 estão
apenas parcialmente atendidos: ainda não há versão de política persistida,
autenticação individual nem registro integrado de confirmação física. O
override usa um token Bearer compartilhado e recebe o identificador de
supervisor informado na requisição; isso demonstra o fluxo, mas não comprova a
identidade do ator. O firmware experimental emite confirmações, e o backend já
possui um adaptador MQTT que valida identificador, sequência, sessão, estado e
bico antes de aceitar o ACK. A integração ao broker ainda não foi ensaiada ao
vivo, e o ACK não constitui uma entidade persistida independente. Mesmo quando
coerente, ele confirma o comando e o estado lógico da saída GPIO, não a posição
mecânica da trava ou do bico.

Os requisitos serão validados com o orientador e, depois das autorizações
necessárias, com representantes dos papéis operacionais. Alterações serão
rastreáveis até a motivação, o cenário e o teste correspondente.

\section{Arquitetura}
\label{sec:arquitetura}

A arquitetura separa interface, decisão, integração externa e acionamento. A
PWA não acessa SENATRAN ou dispositivo diretamente. O backend consulta um
\textit{VehicleDataProvider}, entrega apenas categorias normalizadas ao motor de
política e, se autorizado, envia ao \textit{DispenserGateway} um comando sem
placa ou conteúdo da resposta veicular. A Figura~\ref{fig:arquitetura} apresenta
essa organização.

\begin{figure}[H]
  \centering
  \noindent\resizebox{0.98\textwidth}{!}{\begin{tikzpicture}[
    node distance=6mm and 8mm,
    bloco/.style={draw, rounded corners, align=center, minimum height=9mm,
      text width=25mm, fill=black!3, font=\small},
    externo/.style={bloco, dashed},
    seta/.style={-{Latex[length=2mm]}, thick}
  ]
    \node[bloco] (ui) {PWA\\frentista/\\supervisor};
    \node[bloco, right=of ui] (api) {API e controle\\de autorização};
    \node[bloco, right=of api] (policy) {Política de\\compatibilidade};
    \node[bloco, below=of api] (audit) {Auditoria\\minimizada};
    \node[externo, above=of policy] (provider) {Provedor veicular\\simulado/\\autorizado};
    \node[bloco, right=of policy] (gateway) {Gateway do\\dispensador};
    \node[externo, right=of gateway] (device) {Simulador ou\\ESP32 de bancada};

    \draw[seta] (ui) -- (api);
    \draw[seta] (api) -- (provider);
    \draw[seta] (provider) -- (policy);
    \draw[seta] (api) -- (policy);
    \draw[seta] (policy) -- (gateway);
    \draw[seta] (gateway) -- (device);
    \draw[seta] (api) -- (audit);
    \draw[seta] (policy) -- (audit);
  \end{tikzpicture}}
  \caption{Arquitetura lógica-alvo do protótipo}
  \label{fig:arquitetura}
  \fontepropria
\end{figure}

\subsection{Interfaces principais}

A operação principal já é exposta pelo backend simulado como
\texttt{POST} \path{/api/v1/authorizations}, recebendo placa, combustível solicitado,
identificador do frentista, bomba e um UUID \path{client_request_id}. Repetir a
mesma chave com o mesmo corpo retorna a decisão persistida sem novo comando; a
reutilização com conteúdo diferente retorna conflito. A resposta tem estado,
identificador da autorização, categorias compatíveis, código de motivo e
expiração. Dados pessoais não são devolvidos à interface.

O override é uma operação separada vinculada a uma autorização em
\path{REVIEW_REQUIRED} causada por incerteza da consulta veicular. O contrato
atual exige token Bearer, identificador declarado de supervisor e justificativa;
autenticação individual permanece requisito não implementado. O contrato-alvo
do dispositivo aceitará somente
\path{LOCK_ALL} ou \path{UNLOCK(nozzle_id,ttl)}, acompanhados de
\path{command_id}, sequência monotônica, sessão, horário de emissão e
expiração, conforme o tipo da operação. Cada comando produzirá confirmação
associada ao mesmo identificador e à mesma sequência.

\subsection{Decisões tecnológicas}

O incremento atual utiliza Python/FastAPI no backend, React na PWA e SQLite para
persistência local minimizada. Um gateway em memória representa o dispensador.
Há firmware experimental que implementa o contrato MQTT para ESP32, mas ele
ainda não está integrado ao backend. PostgreSQL e contêineres permanecem opções
para uma etapa posterior, não componentes já empregados. Essas escolhas são
substituíveis: as regras de compatibilidade permanecem no domínio, e integrações
são isoladas pelas interfaces \path{VehicleDataProvider} e
\path{DispenserGateway}.

Somente o modo sintético está habilitado: a configuração do backend rejeita a
inicialização em modo SENATRAN, e o adaptador correspondente não realiza chamada
externa. Essa barreira evita uso acidental de dados ou credenciais antes da
autorização formal.

O ESP32 é o controlador inicialmente recomendado para a bancada, acionando LEDs
e, após revisão do orientador, servo ou solenoide de baixa tensão. Um Raspberry
Pi poderá atuar como gateway local, mas não como mecanismo único de segurança.
Para campo, a alternativa prioritária a estudar é integração com controlador de
pista ou dispenser certificado. A documentação do Fuel Point PLUS, por exemplo,
descreve identificação veículo-bico, consulta a servidor de autorização e
abertura da bomba somente depois da aprovação \cite{gilbarcoFuelPoint2018}; ela
serve como referência de arquitetura comercial, não como validação do protótipo.

\subsection{Disponibilidade e falhas}

O sistema não adotará \textit{fail-open}. Se o provedor estiver indisponível ou
se a resposta não puder ser mapeada de maneira inequívoca, a decisão automática
será \path{REVIEW_REQUIRED}. A organização deverá definir um procedimento
manual seguro para continuidade operacional; o protótipo apenas registrará o
override autorizado. Perda de comunicação depois da liberação fará o
controlador bloquear localmente após o TTL, independentemente do backend.

Se a PWA perder a resposta de criação, ela não presume bloqueio e não cria outra
solicitação automaticamente: a reconciliação manual repete exatamente a mesma
chave e o mesmo corpo. Se a resposta de override ou conclusão se perder, a
interface consulta o estado pelo identificador já conhecido, sem repetir o
comando físico. \path{PENDING_DEVICE} é tratado como estado físico desconhecido.

\section{Modelagem}
\label{sec:modelagem}

\subsection{Estados da autorização}

A Figura~\ref{fig:estados} representa a máquina de estados conceitual almejada.
A consulta começa em \path{QUERYING}. Produto incompatível conduz a
\path{DENIED}; ausência, ambiguidade ou código desconhecido conduz a
\path{REVIEW_REQUIRED}; e uma decisão compatível primeiro persiste a intenção
\path{PENDING_DEVICE}. Somente o reconhecimento do comando promove a
autorização para \path{AUTHORIZED}. Uma revisão excepcional permitida também
passa por essa intenção, nunca diretamente para autorização. Expiração,
conclusão, substituição ou falha retornam o dispositivo ao bloqueio.

\begin{figure}[H]
  \centering
  \noindent\resizebox{0.98\textwidth}{!}{\begin{tikzpicture}[
    node distance=8mm and 10mm,
    estado/.style={draw, rounded corners, align=center, minimum height=9mm,
      text width=25mm, fill=black!3, font=\small},
    seta/.style={-{Latex[length=1.8mm]}, thick, font=\scriptsize}
  ]
    \node[estado] (locked) {LOCKED};
    \node[estado, right=of locked] (query) {QUERYING};
    \node[estado, right=of query] (pending) {PENDING\_\\DEVICE};
    \node[estado, right=of pending] (authorized) {AUTHORIZED};
    \node[estado, right=of authorized] (completed) {COMPLETED};
    \node[estado, above=of query] (denied) {DENIED};
    \node[estado, below=of query] (review) {REVIEW\_\\REQUIRED};
    \node[estado, below=of authorized] (terminal) {EXPIRED/\\SUPERSEDED};

    \draw[seta] (locked) -- (query);
    \draw[seta] (query) -- node[above]{compatível} (pending);
    \draw[seta] (query) -- node[left]{incompatível} (denied);
    \draw[seta] (query) -- node[left]{incerto} (review);
    \draw[seta] (review) -- node[below,sloped]{override} (pending);
    \draw[seta] (pending) -- node[above]{ACK} (authorized);
    \draw[seta] (authorized) -- (completed);
    \draw[seta] (authorized) -- (terminal);
  \end{tikzpicture}}
  \caption{Máquina de estados conceitual da autorização}
  \label{fig:estados}
  \fontepropria
\end{figure}

Estados terminais, expiração, substituição e falhas conduzem ao estado físico
bloqueado; esses retornos foram omitidos da figura para preservar a
legibilidade. A implementação atual persiste \path{PENDING_DEVICE},
\path{AUTHORIZED}, \path{DENIED}, \path{REVIEW_REQUIRED},
\path{COMPLETED}, \path{EXPIRED} e \path{SUPERSEDED};
\path{QUERYING}, \path{DISPENSING}, \path{LOCKED} e \path{CANCELLED}
ainda não são estados de domínio materializados. Não existe transição de
\path{DENIED}, \path{EXPIRED}, \path{COMPLETED} ou
\path{SUPERSEDED} para \path{AUTHORIZED}. O override parte exclusivamente de
\path{REVIEW_REQUIRED} e somente quando o motivo pertence às incertezas da
consulta; falha de bomba, bico ou dispensador não pode ser contornada.

Antes de enviar \path{UNLOCK}, o backend persiste a intenção e o identificador
do comando. O gateway simulado aplica e confirma a operação de modo síncrono;
então o backend registra o evento de ACK e promove a autorização. O firmware
experimental já produz \path{COMMAND_ACK}, porém o backend ainda não consome
esse evento por MQTT. Se o bloqueio compensatório falhar, a garantia residual é
o TTL local do dispositivo, não um bloqueio imediato absoluto.

\subsection{Tabela de decisão}

\begin{table}[H]
  \centering
  \caption{Política preliminar de decisão}
  \label{tab:decisao}
  \small
  \begin{tabularx}{\textwidth}{@{}L{2.4cm}L{2.7cm}L{2.4cm}Y@{}}
    \toprule
    \textbf{Resposta} & \textbf{Produto solicitado} & \textbf{Decisão} & \textbf{Efeito físico} \\
    \midrule
    Válida e inequívoca & Compatível & Autorizar & Liberar apenas o bico solicitado por TTL. \\
    Válida e inequívoca & Incompatível & Negar & Manter todos bloqueados. \\
    Inválida ou ausente & Qualquer & Revisar & Manter bloqueado; permitir fluxo de supervisor. \\
    Código desconhecido & Qualquer & Revisar & Manter bloqueado e registrar versão da política. \\
    Válida com múltiplos & Um compatível & Autorizar & Liberar somente o produto escolhido, nunca todos. \\
    Conversão não comprovada & Qualquer & Revisar & Não inferir compatibilidade automaticamente. \\
    \bottomrule
  \end{tabularx}
  \\[2pt]\fontepropria
\end{table}

\subsection{Modelo de dados minimizado}

O modelo-alvo prevê \path{Authorization}, \path{PolicyVersion},
\path{DeviceCommand}, \path{DeviceAck}, \path{Override} e
\path{AuditEvent}. Na implementação corrente, o SQLite contém somente tabelas
de autorização e eventos de auditoria; comandos ficam no gateway em memória, e
os campos de override integram a autorização. Ainda não existem entidades
persistidas para versão da política ou confirmação do dispositivo.

A placa existe apenas em memória durante a consulta e não é chave de
relacionamento persistente. O esquema atual guarda decisão, combustível
solicitado, bomba, frentista, modo do provedor, motivo, horários e eventual
override, além de \path{client_request_id} e um fingerprint HMAC para verificar
idempotência. A consulta de compatibilidade usa cache temporário em memória,
indexado por HMAC da placa normalizada, para evitar uma segunda consulta no
mesmo fluxo. A chave HMAC fica fora do repositório; esse recurso reduz exposição,
mas é pseudonimização técnica, não anonimização. Não há colunas para placa,
Renavam, chassi, proprietário ou resposta bruta do provedor. Logs de pesquisa
serão separados dos dados operacionais e conterão somente casos sintéticos.

\pendente{validar o modelo e a enumeração oficial de códigos de combustível com
a documentação técnica liberada no ambiente autorizado.}


% (6) Cronograma com estimativa de término.
\section{Cronograma}
\label{sec:cronograma}

Como a data oficial do Seminário de Andamento ainda não foi registrada neste
repositório, o cronograma usa $S$ para representar a semana da apresentação. As
datas serão convertidas assim que houver comunicação institucional confirmada.

\begin{table}[H]
  \centering
  \caption{Cronograma relativo até o Seminário de Andamento}
  \label{tab:cronograma}
  \small
  \begin{tabularx}{\textwidth}{@{}L{2cm}Y@{}}
    \toprule
    \textbf{Período} & \textbf{Entregas verificáveis} \\
    \midrule
    $S-8$ a $S-7$ & Confirmar normas; completar metadados; congelar protocolo de revisão; preparar instrumentos, autorizações e solicitação de dados. \\
    $S-7$ a $S-5$ & Executar buscas; selecionar estudos; consolidar contextualização, problema, objetivos e fundamentação. \\
    $S-5$ a $S-3$ & Coletar dados autorizados; realizar entrevistas aprovadas; validar requisitos e arquitetura; consolidar os simuladores. \\
    $S-3$ a $S-2$ & Ampliar testes de componentes e integração; calcular primeira análise econômica; registrar limitações e resultados sem extrapolação. \\
    $S-2$ a $S-1$ & Submeter ao orientador; corrigir texto, referências, tabelas e compilação; fechar versão rastreável. \\
    $S-1$ a $S$ & Produzir apresentação modular, revisar o PDF, ensaiar e reservar margem para correções. \\
    \bottomrule
  \end{tabularx}
  \\[2pt]\fontepropria
\end{table}

Após o seminário, as atividades serão: finalizar bancada, concluir integração
autorizada se disponível, ampliar os testes, executar a avaliação com usuários,
fechar a análise econômica, revisar ameaças à validade e redigir implementação,
demonstração, resultados, conclusão e defesa final.

\pendente{inserir data oficial do seminário, duração, formato, limite de páginas
e marcos do orientador; então substituir as semanas relativas por datas.}


\begingroup
\footnotesize
\bibliographystyle{sbc}
\bibliography{referencias}
\endgroup

\end{document}
